% !TEX root = creche.tex
\documentclass[creche.tex]{subfiles}
\begin{document}
\chapter{Continous logic}
\label{measures}

\def\medrel#1{\parbox[t]{5ex}{$\displaystyle\hfil #1$}}
\def\ceq#1#2#3{\parbox[t]{18ex}{$\displaystyle #1$}\medrel{#2}{$\displaystyle #3$}}


\def\Msr{{\rm Msr\kern.0ex}}


\section{Real-valued structures}

Let $L$ be a first-order signature. A \emph{real-valued structure\/} or \emph{model\/} $M$ of signature $L$ consists of\nobreak
\begin{itemize}
\item[1] A set that is the domain (or support) of the structure;
\item[2.] for every  $f\in L_{\rm fun}$, a continuous function $f^M:\ M^{n_f}\to M$;
\item[3.] for every $r\in L_{\rm rel}$, a real valued function $r^M:\ M^{n_r}\to \RR$.
\end{itemize}

Classical structures are identified with real valued structures where relation symbos are interpreted with $\{0,1\}$-valued functions: $0$ for true and $1$ for false (sic).

A set of propositional connectives $C$ is a set of function $\RR^n\to\RR$. We assume that $C$ contains at least the following functions

\begin{itemize}
\item[0.] the numbers $0$, $1$ as zero-ary connectives;
\item[1.] all real numbers as unary connectives;
\item[2.] $+$ and $\cdot$ as binary connectives.
\end{itemize}

Atomic formulas are defined as in first-order logic: they are of the form $r\,t$, where $r$ is a relation symbol and $t$ is a tuple of terms. All other formulas are obtained inductively as follows

\begin{itemize}
\item[1.] if $\phi_1,\dots,\phi_n$ are formulas and $c\in C$, then $c\,\phi_1\dots\phi_n$ is a formula;
\item[2.] if $\phi$ is a formula, $x$ is a variable, then $\inf_x\phi$ is a formula.
\end{itemize}

We write $CL$ for the set of formulas obtained from $L$ and $C$. If $A\subseteq M$ we write $CL(A)$ for the set of formulas with parameters in $A$.

Let $M$ be a model and let $t$ be a tuple of closed term possibly with parameters from $M$. We write $t^M$ for the element obtained interpreting function symbols and parameters in $M$ just as in classical first-order logic. If $\phi\in CL(M)$ is an atomic sentence (i.e.\@ a closed atomic formula) we define $\phi^M$ as follows:

\begin{itemize}
\item[1.] $(r\,t)^M\ =\  r^M\big(t^M\big)$.
\end{itemize}

Then, by induction we define $\phi^M$ for any sentence $\phi\in L(M)$:

\begin{itemize}
\item[1.] if $\phi=c\,\phi_1\,\dots\,\phi_n$ then $\phi^M=c(\phi^M_1,\dots,\phi^M_n)$;
\item[2.] if $\phi=\inf_x\psi$ then $\phi^M=\inf\{\psi[x/a]^M: a\in M\}$.
\end{itemize}


For $\phi, \psi\in CL(M)$, we may write \emph{$M\models\phi=\psi$\/} for $\phi^M=\psi^M$. We also define \emph{$M\models\ \phi\neq\psi$\/} in a similar way. Expressions of the form $\phi=0$ are called \emph{conditions}. Expressions of the form $\phi\neq 0$ are called \emph{co-conditions}.

Let $A\subseteq M$. The \emph{$A$-topology\/} on $M^{|x|}$ is the topology generated by open sets defined by co-conditions


\ceq{\hfill \emph{$\phi(M){\neq}0$}}{=}{\Big\{a\in M^{|x|}\ :\ M\models\phi(a)\neq 0\Big\}}

for some $\phi(x)\in CL(A)$.

\begin{proposition}\label{prop_CL_Tarski-Vaught} Let $N$ be a model and and let $M\subseteq N$ be a substructure. Fix a tuple of variables $x$ of finite length. Then the following are equivalent:
\begin{itemize}
\item[1.] $M\preceq N$;
\item[2.] $M^{|x|}$ is dense in $N^{|x|}$ w.r.t.\@ the $M\mbox{-}$topology on $N^{|x|}$. That is, for every $\phi(x)\in L(M)$, if $\phi(N){\neq}0$ is non empty, then $\phi(b){\neq}0$ for some $b\in M^{|x|}$.
\end{itemize}
\end{proposition}
\begin{proof}
Implication \ssf{1}$\IMP$\ssf{2} is clear. We prove \ssf{2}$\IMP$\ssf{1} by induction on the syntax of $\phi(x)$. We use \ssf{\#} as induction hypothesis. As $M$ is a substructure \ssf{\#} holds for atomic $\phi(x)$. We leave the case of propositional connectives to the reader and consider only the case of the quantifier.

\ceq{\hfill N\,\models\,\inf_x\phi(x){<}0}{\IFF}{} the set $\phi(N)<\epsilon$ is non empty for some $\epsilon>0$

\ceq{\ssf{a.}}{\IFF}{} there is $b\in M^{|x|}$ such that $N\models\phi(b)<\epsilon$

\ceq{\ssf{b.}}{\IFF}{} there is $b\in M^{|x|}$ such that $M\models\phi(b)<\epsilon$

\ceq{}{\IFF}{M\,\models\,\bigwedge^{\phantom{n}}_x\phi(x)<\epsilon.}

Implication $\IMP$ in \textsf{a} uses \ssf{2} and equivalence \textsf{b} uses the induction hypothesis.
\end{proof}







Equations of the form $\phi(x)=0$ are called \emph{conditions}. Observe that $\phi(x)=0$ is equivalent to $|\phi(x)|\le0$ and $\phi(x)\le0$ is equivalent to $\phi(x)\vee0=0$. So, when convenient we may use weak inequalities to denote conditions. We write \emph{$\phi(M)=0$} and \emph{$\phi(M)\le0$} for the corresponding sets of solutions in $M$. As $\phi(x)$ ranges over $L(M)$ these are the zero sets of the limit topology.

A set of conditions is called a \emph{type\/} or, when there are no free variables, a \emph{theory}. For any $p(x)\subseteq L(M)$ we write

\ceq{\hfill\emph{$p(x)\le0$}}{\deq}{\bigg\{\phi(x)\le0 \quad :\quad  \phi(x)\in p\bigg\}}

We write \emph{$M\models p(a)\le0$} if $M\models\phi(a)\le0$ for every $\phi(x)\in p$. We write $p(M)\le0$ for the set of solutions of $p(x)\le0$. This is a closed set in the limit topology. All closed sets of the limit topology have this form. 


We write $p^M$ for the function $p^M:M^{|x|}\to\RR\cup\{\infty\}$ that maps $a\mapsto\sup\big\{\phi(a)^M\ :\ \phi(x)\in p\big\}$. Then it is immediate that

\ceq{\hfill M\models p(a)\le 0}{\IFF}{p^M(a)\le 0}.

We will not not use expressions like \emph{$p(x)\le0$} or \emph{$p(x)=0$} because they do not satisfy the equivalence above.



Let $\Delta\subseteq L(M)$ be a set of formulas closed under Riesz connectives. For $a\in M^{|x|}$ we write \emph{$p(x)=\Deltatp_M(a)$\/} to say that 

\ceq{\hfill p(x)}{=}{\Big\{\phi(x)\in\Delta\ :\ M\models\phi(a)\le 0\Big\}.}

The subscript $M$ will be omitted when clear from the context and when $\Delta=L(A)$ we will write \emph{$p(x)=\tp_M(a/A)$}.  Note that for every $\phi(x)$ either $\phi(x)$ or $\neg\phi(x)$ belong to $p(x)=\tp_M(a/A)$. Hence, $M\models 0\le p^M(b)$ for any $b\in M^{|x|}$. But this may not be true in general. In fact, for some $p(x)\subseteq L(A)$, the there may be no type equivalent to 

\ceq{~}{~}{\Big\{\phi(x)\in\Delta\ :\ \phi(x)\in p,\ \ M\models0\le \phi(a)\Big\}.}

Indeed, we avoid espressiones as $p(x)\ge0$, $p(x)=0$ as for these an equivalence as in $\#$ may not hold.


\end{document}
